% ============================================================================
% SECCIÓN 1: INTRODUCCIÓN
% ============================================================================

\section{Introducción}

\subsection{Contexto del Proyecto}

El presente informe corresponde al proyecto final de la asignatura \textbf{EYP2605 - Matemáticas Actuariales}, desarrollado durante noviembre de 2025. El objetivo principal es construir una tabla de mortalidad utilizando datos reales de la población asegurada, aplicando técnicas de graduación y extrapolación estudiadas durante el curso.

\subsection{Datos y Especificaciones del Proyecto}

Los parámetros asignados para este proyecto son:

\begin{hallazgobox}
\textbf{Especificaciones del Proyecto:}
\begin{itemize}
    \item \textbf{Tabla Asignada:} INVALIDEZ 2014 - Masculino
    \item \textbf{Origen:} 2 (Origen = 2)
    \item \textbf{Sexo:} Masculino (TM\_SEXO = M)
    \item \textbf{Período de Análisis:} 01/01/2000 -- 31/12/2012
    \item \textbf{Tasa de Interés Técnico:} 5\% anual
\end{itemize}
\end{hallazgobox}

\subsection{Objetivos del Informe}

Este proyecto se estructura en tres componentes principales, cada uno con su ponderación respectiva:

\begin{enumerate}
    \item \textbf{Construcción de Tabla de Mortalidad (60\%):}
    \begin{itemize}
        \item Obtención del cuerpo de la tabla mediante graduación
        \item Análisis del período 2000-2012
        \item Graduación con evaluación de bondad de ajuste
        \item Extrapolación para completar cabeza y cola
        \item Construcción de tabla completa con esperanza de vida
    \end{itemize}
    
    \item \textbf{Valores de Conmutación (10\%):}
    \begin{itemize}
        \item Cálculo de funciones de conmutación para la tabla obtenida
        \item Cálculo de funciones de conmutación para la tabla de referencia
        \item Uso de tasa de interés técnico del 5\%
    \end{itemize}
    
    \item \textbf{Cálculo de Valores Actuales (30\%):}
    \begin{itemize}
        \item Dotal puro
        \item Renta incierta vitalicia con pago vencido
        \item Renta incierta vitalicia diferida con pago anticipado
        \item Renta incierta temporal con pago vencido
        \item Renta incierta temporal diferida con pago anticipado
        \item Comparación entre tabla obtenida y tabla de referencia
    \end{itemize}
\end{enumerate}

\subsection{Estructura del Informe}

El documento se organiza de la siguiente manera:

\begin{itemize}
    \item \textbf{Sección 2:} Construcción de la tabla de mortalidad, incluyendo análisis de datos, graduación, extrapolación y tabla completa.
    \item \textbf{Sección 3:} Cálculo de valores de conmutación para ambas tablas.
    \item \textbf{Sección 4:} Cálculo de valores actuales y comparación de resultados.
    \item \textbf{Sección 5:} Conclusiones y análisis de los hallazgos.
\end{itemize}

\subsection{Metodología}

Para el desarrollo de este proyecto se utilizó el lenguaje de programación \textbf{R}, implementando un proceso estructurado en las siguientes etapas:

\begin{enumerate}
    \item \textbf{Preparación de datos:} Filtrado por Origen=2 y Sexo=M, conversión de fechas y limpieza de datos.
    
    \item \textbf{Cálculo de exposición:} Método actuarial con funciones \texttt{pmin/pmax} vectorizadas para cálculo eficiente de la exposición al riesgo.
    
    \item \textbf{Suavizamiento:} Método Whittaker-Henderson aplicado en el rango 20-84 años.
    
    \item \textbf{Extrapolación edades jóvenes:} Método de Oppermann para edades 0-19 años.
    
    \item \textbf{Extrapolación edades avanzadas:} Método de Heligman-Pollard para edades 85-109 años.
    
    \item \textbf{Construcción de tabla de vida:} Cálculo de funciones biométricas ($q_x$, $l_x$, $d_x$, $e_x$) y funciones de conmutación ($D_x$, $N_x$, $S_x$, $C_x$, $M_x$, $R_x$).
    
    \item \textbf{Pruebas de bondad de ajuste:} Chi-cuadrado, Kolmogorov-Smirnov, Test de Rachas y Test de Signos.
    
    \item \textbf{Cálculo de valores actuales:} Para diversos productos de seguros y rentas.
\end{enumerate}

\textbf{Librerías de R utilizadas:}

\begin{itemize}
    \item \texttt{tidyverse}, \texttt{lubridate} — Manipulación y análisis de datos
    \item \texttt{WH}, \texttt{MortalityTables} — Suavizamiento Whittaker-Henderson
    \item \texttt{MortalityLaws} — Métodos de extrapolación
    \item \texttt{randtests}, \texttt{BSDA}, \texttt{tseries} — Tests estadísticos de bondad de ajuste
\end{itemize}

El script principal se encuentra en: \texttt{C\_Optimizadosycompletos.R}
